% =============================================================================
% Worked solutions for Simulados/mock-midterm-topics-1-2b.pdf (Topics 1, 2A, 2B).
% Statements: mock-midterm-topics-1-2b-statements.tex (shared with the exam paper).
% Compile: python -I Simulados/mock-midterm-topics-1-2b-solutions_check.py (must print PASS),
%          then pdflatex (twice); check: pdffonts (Type 1 only, uni = yes).
% =============================================================================
\documentclass[11pt,a4paper]{article}
\usepackage[utf8]{inputenc}
\usepackage[T1]{fontenc}
\usepackage{lmodern}   % Type 1 vector fonts (replaces PK bitmaps)
\usepackage{cmap}      % ToUnicode CMap -> copyable text and accents
\usepackage[english]{babel}
\usepackage{amsmath,amssymb,amsfonts}
\usepackage[a4paper,margin=2.2cm]{geometry}
\usepackage{booktabs}
\usepackage{enumitem}
\usepackage{fancyhdr}
\usepackage{xcolor}
\usepackage{microtype}

\pagestyle{fancy}\fancyhf{}
\lhead{\footnotesize Finan\c{c}as II -- Corporativas (MPE Insper)}
\rhead{\footnotesize Mock midterm, Topics 1--2B --- Worked solutions}
\cfoot{\footnotesize\thepage}
\renewcommand{\headrulewidth}{0.3pt}
\setlength{\headheight}{14pt}

\makeatletter
\renewcommand{\@seccntformat}[1]{\csname the#1\endcsname.\quad}
\makeatother
\numberwithin{equation}{section}
\setlist{topsep=2pt,itemsep=1pt}

\newcommand{\srcline}[1]{\par\noindent{\small\color{black!65}\textbf{Source:} #1}\par\smallskip}
\newcommand{\statement}[1]{\par\smallskip\noindent\colorbox{black!5}{\parbox{\dimexpr\linewidth-2\fboxsep}{\small #1}}\par\medskip}
\newcommand{\step}[2]{\par\medskip\noindent\textbf{Step #1 --- #2.}~}
\newenvironment{notation}{\par\noindent\textbf{Notation.}\begin{itemize}[leftmargin=1.4em]}{\end{itemize}}
\newcommand{\intuition}[1]{\par\medskip\noindent\textbf{Intuition.}~#1}
\newcommand{\refline}[1]{\par\smallskip\noindent{\small\color{black!70}\textbf{Ref:} #1}}
\newcommand{\rubric}[1]{\par\smallskip\noindent{\small\textbf{Marking guide (10 marks).}~#1}}
\newcommand{\answer}[1]{\par\medskip\noindent\fbox{\parbox{\dimexpr\linewidth-2\fboxsep-2\fboxrule}{\textbf{Answer.}~#1}}\par}

% =============================================================================
% Statements of the mock midterm (Topics 1, 2A, 2B), written once and \input by
%   mock-midterm-topics-1-2b.tex            (exam paper)
%   mock-midterm-topics-1-2b-solutions.tex  (worked solutions)
% Every statement is a real exercise copied verbatim (rules/exercise-lists.md):
% source line first, then the full text. Allowed edits only: math in LaTeX, tables
% rebuilt, line breaks undone, currency written as \$.
% Needs: \srcline{...} defined by the including document.
% =============================================================================

% ---------------------------------------------------------------- Question 1.1
\newcommand{\QoneA}{%
\srcline{Berk, DeMarzo \& Harford, \emph{Fundamentals of Corporate Finance} 2e, Ch.~2, Exercise 2.24,
p.~55 (PDF 86) $\cdot$ Answer in book: no}

\noindent In January 2009, American Airlines (AMR) had a market capitalization of \$1.7 billion, debt of
\$11.1 billion, and cash of \$4.6 billion. American Airlines had revenues of \$23.8 billion. British
Airways (BABWF) had a market capitalization of \$2.2 billion, debt of \$4.7 billion, cash of \$2.6
billion, and revenues of \$13.1 billion.
\begin{enumerate}[label=\alph*.,leftmargin=2em]
  \item Compare the market capitalization-to-revenue ratio (also called the price-to-sales ratio) for
        American Airlines and British Airways.
  \item Compare the enterprise value-to-revenue ratio for American Airlines and British Airways.
  \item Which of these comparisons is more meaningful? Explain.
\end{enumerate}}

% ---------------------------------------------------------------- Question 1.2
\newcommand{\QoneB}{%
\srcline{Berk, DeMarzo \& Harford, \emph{Fundamentals of Corporate Finance} 2e, Ch.~9, Exercise 9.13,
p.~275 (PDF 306) $\cdot$ Answer in book: no}

\noindent Elmdale Enterprises is deciding whether to expand its production facilities. Although long-term
cash flows are difficult to estimate, management has projected the following cash flows for the first
two years (in millions of dollars):
\begin{center}
\begin{tabular}{lrr}
  \toprule
  Year & 1 & 2 \\ \midrule
  Revenues & 125 & 160 \\
  Operating Expenses (other than depreciation) & 40 & 60 \\
  Depreciation & 25 & 36 \\
  Increase in Net Working Capital & 2 & 8 \\
  Capital Expenditures & 30 & 40 \\
  Marginal Corporate Tax Rate & 35\% & 35\% \\ \bottomrule
\end{tabular}
\end{center}
\begin{enumerate}[label=\alph*.,leftmargin=2em]
  \item What are the incremental earnings for this project for years 1 and 2?
  \item What are the free cash flows for this project for the first two years?
\end{enumerate}}

% ---------------------------------------------------------------- Question 1.3
\newcommand{\QoneC}{%
\srcline{Berk, DeMarzo \& Harford, \emph{Fundamentals of Corporate Finance} 2e, Ch.~16, Exercises
16.19--16.21, pp.~493--494 (PDF 524--525) $\cdot$ Answer in book: no}

\noindent\textbf{19.} Kohwe Corporation plans to issue equity to raise \$50 million to finance a new
investment. After making the investment, Kohwe expects to earn free cash flows of \$10 million each year.
Kohwe currently has 5 million shares outstanding, and it has no other assets or opportunities. Suppose
the appropriate discount rate for Kohwe's future free cash flows is 8\%, and the only capital market
imperfections are corporate taxes and financial distress costs.
\begin{enumerate}[label=\alph*.,leftmargin=2em]
  \item What is the NPV of Kohwe's investment?
  \item What is Kohwe's share price today?
\end{enumerate}
\noindent\textbf{20.} Suppose Kohwe borrows the \$50 million instead. The firm will pay interest only on
this loan each year, and it will maintain an outstanding balance of \$50 million on the loan. Suppose
that Kohwe's corporate tax rate is 40\%, and expected free cash flows are still \$10 million each year.
What is Kohwe's share price today if the investment is financed with debt?

\smallskip\noindent\textbf{21.} Now suppose that with leverage, Kohwe's expected free cash flows will
decline to \$9 million per year due to reduced sales and other financial distress costs. Assume that the
appropriate discount rate for Kohwe's future free cash flows is still 8\%. What is Kohwe's share price
today given the financial distress costs of leverage?}

% ---------------------------------------------------------------- Question 1.4
\newcommand{\QoneD}{%
\srcline{Berk, DeMarzo \& Harford, \emph{Fundamentals of Corporate Finance} 2e, Ch.~17, Exercises
17.11--17.13, p.~528 (PDF 559) $\cdot$ Answer in book: no. Problems 12 and 13 say ``Redo Problem 11'',
so Problem 11 is part of this block.}

\noindent\textbf{11.} Assume perfect capital markets. Kay Industries currently has \$100 million invested
in short-term Treasury securities paying 7\%, and it pays out the interest payments on these securities
as a dividend. The board is considering selling the Treasury securities and paying out the proceeds as a
one-time dividend payment.
\begin{enumerate}[label=\alph*.,leftmargin=2em]
  \item If the board went ahead with this plan, what would happen to the value of Kay stock upon the
        announcement of a change in policy?
  \item What would happen to the value of Kay stock on the ex-dividend date of the one-time dividend?
  \item Given these price reactions, will this decision benefit investors?
\end{enumerate}
\noindent\textbf{12.} Redo Problem 11, but assume that Kay must pay a corporate tax rate of 35\%, and that
investors pay no taxes.

\smallskip\noindent\textbf{13.} Redo Problem 11, but assume that investors pay a 15\% tax on dividends but
no capital gains taxes, and that Kay does not pay corporate taxes.}

% ---------------------------------------------------------------- Question 2 (previous-year midterm)
\newcommand{\QtwoIntro}{%
\srcline{Finance II -- Corporate Finance, \emph{Midterm Exam (Previous Year)}, Question 2, pp.~1--2
(\texttt{Prova/q2.pdf}) $\cdot$ Answer in source: no}

\noindent Company A has no assets and is considering investing in a new project. The expected
\emph{after-tax} unlevered free cash flows are given below:
\begin{center}
\begin{tabular}{lrrrrrr}
  \toprule
  Year & 0 & 1 & 2 & 3 & 4 & 5 \\ \midrule
  FCF & $-750$ & 250 & 250 & 250 & 250 & 250 \\ \bottomrule
\end{tabular}
\end{center}
As the project is new, owners of company A will look at comparable firms to find an appropriate discount
rate. They identify a comparable firm whose assets (\emph{excluding the market value of any tax benefits
from debt}) have the same market or systematic risk as the project that company A is considering
implementing. The equity of the comparable firm trades in financial markets. The market value of the
assets (\emph{excluding the market value of any tax benefits from debt}) of the comparable firm is
denoted by $V_U$ (as if the firm had no debt). If the comparable firm had no leverage, the market value
of its assets would be equal to $V_U$, that is, $A=V_U$. However, as the comparable firm maintains a
constant debt-to-equity ratio of 0.1, the market value of its assets equals $V_U$ plus the present value
of interest tax shields, that is, $A=V_U+PV(ITS)$. The debt of the comparable firm is risk-free and its
equity has a beta of $\beta_E=1.65$. The expected excess market return, $\mathbb{E}[\tilde r_m]-r_f$, is
10\% and the risk free rate is $r_f=5\%$. The assumptions of the CAPM and the Modigliani-Miller model hold
(including competitive financial markets). The corporate tax rate is 30\%.}

\newcommand{\QtwoA}{%
\textbf{(a)} (\emph{10 marks}) What is the beta of the total assets ($\beta_A$), assets excluding tax
benefits of debt ($\beta_U$), debt ($\beta_D$), and tax shields ($\beta_{ITS}$) of the comparable firm?
What is the expected return on the assets ($r_A$), assets excluding tax benefits ($r_U$), debt ($r_D$),
tax shields ($r_{ITS}$), and equity ($r_E$) of the comparable firm?\\
\emph{Hint: we have that}
\[ r_A=\frac DA r_D+\frac EA r_E=\frac{V_U}{A}r_U+\frac{PV(ITS)}{A}r_{ITS} \]
\emph{and}
\[ \beta_A=\frac DA\beta_D+\frac EA\beta_E=\frac{V_U}{A}\beta_U+\frac{PV(ITS)}{A}\beta_{ITS} \]
\emph{but we made an assumption about the capital structure of the comparable firm -- constant
debt-to-equity ratio -- that simplifies this formula.}}

\newcommand{\QtwoB}{%
\textbf{(b)} (\emph{10 marks}) Suppose the owners of company A decide to finance the investment using
equity only. What is the expected return on the assets ($r_A$) and equity ($r_E$) of company A?
Calculate the NPV of the project. Assuming the project is undertaken, what is the market value of the
assets of the firm ($V_U$)? Are outside investors willing to provide financing? If yes, what share of the
firm they buy in exchange for providing 750?}

\newcommand{\QtwoC}{%
\textbf{(c)} (\emph{10 marks}) Suppose the owners of company A decide to maintain a constant
debt-to-equity ratio of 0.5. Debt is still risk-free. What is the required return on equity, $r_E$? Is
it smaller, equal, or larger than the expected return on equity found in (b)? Explain.}

\newcommand{\QtwoD}{%
\textbf{(d)} (\emph{10 marks}) Keeping the assumptions of part (c) -- constant debt-to-equity ratio of 0.5
and risk-free debt --, what would be the NPV of the project? Explain the reason for your findings. Are
investors willing to provide financing?}

\newcommand{\QtwoE}{%
\textbf{(e)} (\emph{10 marks}) Keeping the assumptions of part (c) -- constant debt-to-equity ratio of 0.5
and risk-free debt --, what is the market value of debt (D), equity (E), and the total assets of the firm
at time 0? What is the share of equity that original owners retain and its market value? The value of
the position of original owners is smaller, equal, or larger than the NPV of the project? Explain.}

\newcommand{\QtwoF}{%
\textbf{(f)} (\emph{10 marks}) Keeping the assumptions of part (c) --constant debt-to-equity ratio of 0.5
and risk-free debt --, what is the market value of debt at $t=1$, $t=2$, $t=3$, $t=4$, and $t=5$? Use
these values to compute the stream of interest tax shields and its present value, $PV(ITS)$. What is the
relation between $PV(ITS)$, the value of the assets of the levered firm and results of part (b)?\\
\emph{Hint: Interest tax shields at t=1 are equal to $D_0\times r_D\times\tau$, where $D_0$ is market
value of debt at $t=0$; interest tax shields at t=2 are equal to $D_1\times r_D\times\tau$, where $D_1$
is market value of debt at $t=1$; and so on.}}


\begin{document}

\begin{center}
  {\Large\bfseries Finance II --- Mock Midterm Exam (Topics 1, 2A, 2B)}\\[4pt]
  {\large Worked solutions}
\end{center}

{\sloppy\noindent Exam paper: \texttt{Simulados/mock-midterm-topics-1-2b.pdf}. Each question's full statement
(grey box, with its source) comes first, then the solution step by step, a boxed answer and a marking
guide. Money is in the units of each statement (\$ billion in 1.1, \$ million in 1.2--1.4). Rates are
decimals inside equations. Every number is reproduced by
\texttt{Simulados/mock-midterm-topics-1-2b-solutions\_check.py}. References: lecture slides
\texttt{Aula\_Topic1}, \texttt{Aula\_Topic2A}, \texttt{Aula\_Topic2B} and the formula sheet
\texttt{Notas/formula-sheet-topics-1-2b.pdf}.\par}

% =============================================================================
\section{Question 1.1 --- American Airlines vs.\ British Airways: which multiple?}

\statement{\QoneA}

\begin{notation}
  \item $MC$: market capitalization (market value of equity); $D$: debt; $C$: cash; $R$: revenues.
  \item $EV=MC+D-C$: enterprise value, the market value of the operating business (net of cash),
        claimed jointly by equity and net debt.
\end{notation}

\subsection*{(a) Price-to-sales}
\[
  \frac{MC}{R}\Big|_{AMR}=\frac{1.7}{23.8}=0.0714,\qquad
  \frac{MC}{R}\Big|_{BA}=\frac{2.2}{13.1}=0.1679 .
\]
On this measure British Airways looks about $0.1679/0.0714\approx2.35$ times as expensive per dollar of
sales.

\subsection*{(b) Enterprise value-to-revenue}
\step{1}{Enterprise values}
$EV_{AMR}=1.7+11.1-4.6=8.2$ and $EV_{BA}=2.2+4.7-2.6=4.3$ (\$ billion).

\step{2}{Ratios}
\[
  \frac{EV}{R}\Big|_{AMR}=\frac{8.2}{23.8}=0.3445,\qquad
  \frac{EV}{R}\Big|_{BA}=\frac{4.3}{13.1}=0.3282 .
\]
Once net debt is added back, the two airlines are valued almost the same per dollar of revenue.

\subsection*{(c) Which comparison is meaningful}
EV/revenue. Revenue is produced by the operating assets, which are financed by \emph{both} debt and
equity, so the matching numerator is the value of the whole enterprise. Market capitalization is only
the equity slice. Its size relative to revenue depends on leverage: American's net debt is
$11.1-4.6=6.5$ against $2.1$ for British Airways, so its equity is a thin residual claim, and the
price-to-sales ratio makes it look ``cheap'' only because it is more levered.

\answer{(a) P/S: AMR $0.071$, BA $0.168$. (b) EV/R: AMR $0.345$, BA $0.328$. (c) EV/revenue is the
meaningful comparison: numerator and denominator both refer to the whole firm, so it is not distorted by
the very different leverage of the two airlines.}

\intuition{This is MM read on a balance sheet: the left-hand side (operating value) is the same
regardless of how the right-hand side is split. A multiple that mixes a firm-level flow (sales, EBITDA)
with an equity-level value mixes the two sides.}
\refline{Topic 1 slides (market-value balance sheet, enterprise value); formula sheet, Topic 1 block.}
\rubric{3 for each ratio pair in (a) and (b) with the EV formula stated; 4 for (c): EV/R chosen,
justified by matching claims (firm-level flow with firm-level value) and the leverage distortion.}

% =============================================================================
\section{Question 1.2 --- Elmdale Enterprises: incremental earnings and free cash flow}

\statement{\QoneB}

\begin{notation}
  \item $EBIT_t=\text{Revenues}_t-\text{OpEx}_t-\text{Dep}_t$; $\tau=0.35$.
  \item Incremental (unlevered) earnings $=EBIT_t(1-\tau)$: no interest, because financing is kept out
        of project cash flows.
  \item $FCF_t=EBIT_t(1-\tau)+\text{Dep}_t-\text{CapEx}_t-\Delta NWC_t$.
\end{notation}

\subsection*{(a) Incremental earnings}
\begin{center}
\begin{tabular}{lrr}
  \toprule
  \$ million & Year 1 & Year 2 \\ \midrule
  Revenues & 125 & 160 \\
  $-$ Operating expenses & $-40$ & $-60$ \\
  $-$ Depreciation & $-25$ & $-36$ \\
  $=$ EBIT & 60 & 64 \\
  $-$ Taxes (35\%) & $-21$ & $-22.4$ \\
  $=$ Incremental earnings & \textbf{39} & \textbf{41.6} \\ \bottomrule
\end{tabular}
\end{center}

\subsection*{(b) Free cash flows}
\step{1}{Bridge from earnings}
Add back depreciation (a non-cash charge), then subtract the cash actually spent on capital and working
capital:
\[
  FCF_1=39+25-30-2=32,\qquad FCF_2=41.6+36-40-8=29.6 .
\]
\step{2}{Cross-check with the tax-shield form}
$FCF_t=(\text{Rev}-\text{OpEx})(1-\tau)+\tau\,\text{Dep}-\text{CapEx}-\Delta NWC$:
year 1, $85(0.65)+8.75-30-2=32$; year 2, $100(0.65)+12.6-40-8=29.6$. The terms $\tau\,\text{Dep}$
($8.75$ and $12.6$) are the depreciation tax shields: depreciation matters for cash only through taxes.

\answer{(a) Incremental earnings: \$39 million (year 1) and \$41.6 million (year 2).
(b) Free cash flows: \$32 million (year 1) and \$29.6 million (year 2).}

\intuition{Earnings rise from year 1 to year 2 while free cash flow falls, because the project absorbs
more capital expenditure and working capital in year 2. Investors are paid in cash, so the NPV
discounts FCF, not earnings (PS1, Exercise 1a).}
\refline{Topic 1 slides (NI $\to$ FCF bridge, depreciation tax shield); Topic 2A slides (incremental
cash flows); Berk Ch.~9.}
\rubric{4 for (a), with the EBIT line and no interest deducted; 6 for (b): add back depreciation (2),
subtract CapEx and $\Delta NWC$ with correct signs (3), both years correct (1).}

% =============================================================================
\section{Question 1.3 --- Kohwe Corporation: equity, debt, and financial distress}

\statement{\QoneC}

\begin{notation}
  \item $r=8\%$: discount rate for the free cash flows (unlevered cost of capital); $N_0=5$ million
        shares; $I=50$: investment; $\tau_c=40\%$.
  \item $V_U$: value of the unlevered firm; $V_L=V_U+PV(ITS)$; for permanent debt $D$ whose shields are
        as risky as the debt, $PV(ITS)=\tau_c D$.
\end{notation}

\subsection*{Problem 19 --- equity financing}
\step{1}{NPV}
The FCF are a perpetuity: $V_U=10/0.08=125$, so
\[ NPV=125-50=75. \]
\step{2}{Share price}
In a competitive market new shareholders pay a fair price, so the old shareholders capture the NPV.
Before the issue, existing equity is worth $125-50=75$, so $P=75/5=\$15$. Check: Kohwe sells
$50/15=3.3333$ million shares, and after the investment $8.3333$ million shares $\times\$15=125=V_U$.

\subsection*{Problem 20 --- debt financing}
Permanent debt of 50 at a 40\% tax rate gives $PV(ITS)=0.40\times50=20$. Then $V_L=125+20=145$, equity
$=145-50=95$, and
\[ P=\frac{95}{5}=\$19 . \]
The \$4 gain per share over equity financing is the tax shield ($20/5$), which goes to the existing
shareholders.

\subsection*{Problem 21 --- debt financing with distress costs}
With leverage the expected FCF fall to 9, so $V_U=9/0.08=112.5$; the tax shield is unchanged at 20, so
$V_L=132.5$, equity $=132.5-50=82.5$, and
\[ P=\frac{82.5}{5}=\$16.50 . \]
Equivalently, $V_L=V_U+PV(ITS)-PV(\text{distress})=125+20-12.5$, with $PV(\text{distress})=1/0.08=12.5$.

\answer{19: NPV $=\$75$ million, price $\$15$. 20: $\$19$. 21: $\$16.50$; debt still beats equity
($16.50>15$) because the tax shield (20) exceeds the present value of distress costs (12.5).}

\intuition{This is the static trade-off in three lines: the tax shield raises value, expected distress
costs lower it, and leverage pays while the first exceeds the second. The price moves at once because
all the value change is capitalised into the old shares.}
\refline{Topic 2B slides (ITS, $V_L=V_U+PV(ITS)-PV(\text{distress})$); formula sheet, debt policies table.}
\rubric{3 for 19 (NPV and price, with the reason old shareholders capture the NPV); 3 for 20
($PV(ITS)=\tau_cD$); 4 for 21 (lower $V_U$, unchanged shield, price, and the trade-off reading).}

% =============================================================================
\section{Question 1.4 --- Kay Industries: paying out cash under three tax regimes}

\statement{\QoneD}

\begin{notation}
  \item $C=100$: Treasuries held; $r_f=7\%$; $\tau_c$: corporate tax rate; $\tau_d$: personal tax on
        dividends; capital gains untaxed.
  \item Value created by the policy change $=$ (after-tax value to investors of the one-time payout) $-$
        (after-tax value to investors of keeping the Treasuries and paying out the interest).
\end{notation}

\subsection*{Problem 11 --- perfect capital markets}
(a) No change. Selling Treasuries at their fair price is a zero-NPV transaction; the interest of 7 a
year is worth $7/0.07=100$, exactly the cash that would be paid out. (b) On the ex-dividend date the
stock's total value falls by \$100 million (per share by $100/N$), but shareholders have received \$100
million in cash. (c) No: shareholder wealth is unchanged; a shareholder who wanted the old policy can
buy Treasuries with the dividend (homemade payout policy).

\subsection*{Problem 12 --- corporate tax 35\%, no investor taxes}
Inside the firm the interest is taxed: the after-tax interest paid out is $7(1-0.35)=4.55$ a year,
worth $4.55/0.07=65$ to investors, who can earn 7\% untaxed by themselves. Paid out now, the cash is
worth 100. (a) On the announcement the stock rises by $100-65=\$35$ million, which is the present value
of the corporate taxes on the interest that would no longer be paid ($0.35\times100$). (b) On the
ex-dividend date it falls by the \$100 million dividend. (c) Yes: investors gain \$35 million.

\subsection*{Problem 13 --- 15\% dividend tax, no capital gains tax, no corporate tax}
Status quo: dividends of 7 a year are worth $7(1-0.15)=5.95$ after tax, or $5.95/0.07=85$. One-time
payout: $100(1-0.15)=85$ after tax. (a) No change on the announcement, since $85=85$: the dividend tax is
paid either way, now or year by year. (b) On the ex-dividend date the stock falls by the after-tax value
of the dividend, \$85 million (from $P_{cum}-P_{ex}=Div(1-\tau_d)$ with $\tau_g=0$). (c) No: investors
are indifferent.

\answer{11: (a) no change; (b) falls by \$100 million; (c) no benefit. 12: (a) rises by \$35 million;
(b) falls by \$100 million; (c) yes, $+\$35$ million. 13: (a) no change; (b) falls by \$85 million;
(c) no benefit.}

\intuition{Retaining cash is a financing decision. It is irrelevant in a perfect market, costly when
the firm pays corporate tax on its interest income (corporate tax is the friction that makes excess
cash expensive), and neutral when only dividends are taxed and the tax cannot be avoided by deferring.}
\refline{Topic 2B slides (payout irrelevance, taxes on payout); Berk Ch.~17 (retention of cash, the
effective dividend tax rate).}
\rubric{3 for 11 (zero NPV, ex-dividend drop, homemade payout); 4 for 12 (after-tax interest 4.55, PV 65,
gain 35); 3 for 13 (85 vs.\ 85, ex-dividend drop of 85).}

% =============================================================================
\section{Question 2 --- Previous-year midterm: comparables, constant leverage and tax shields}

\statement{\QtwoIntro}

\begin{notation}
  \item $A=V_U+PV(ITS)=D+E$: market-value balance sheet of the comparable firm (and later of company A).
  \item $\beta_X$, $r_X$: beta and expected return of claim $X\in\{A,U,D,E,ITS\}$; CAPM:
        $r_X=r_f+\beta_X(\mathbb{E}[\tilde r_m]-r_f)=0.05+0.10\,\beta_X$.
  \item $\tau=0.30$; risk-free debt: $\beta_D=0$, $r_D=r_f=5\%$.
  \item Constant $D/E$ (the debt is rebalanced as the firm's value moves) $\Rightarrow$ the tax shields
        are as risky as the unlevered assets: $\beta_{ITS}=\beta_U$, $r_{ITS}=r_U$ (formula sheet,
        ``continuous target $D/V$'').
\end{notation}

% -----------------------------------------------------------------------------
\subsection*{(a) Betas and expected returns of the comparable firm}
\statement{\QtwoA}

\step{1}{Why constant leverage simplifies the hint}
With $D/E$ held constant, debt is always a fixed fraction of firm value, so next year's debt, and hence
next year's tax shield, moves one-for-one with the value of the assets. The shields therefore inherit
the systematic risk of $U$: $\beta_{ITS}=\beta_U$. The right-hand side of the hint becomes
$\frac{V_U}{A}\beta_U+\frac{PV(ITS)}{A}\beta_U=\beta_U$, so $\beta_A=\beta_U$ (and $r_A=r_U$).

\step{2}{Asset beta from the equity beta}
$D/E=0.1\Rightarrow D/A=0.1/1.1=0.0909$ and $E/A=1/1.1=0.9091$. With $\beta_D=0$:
\begin{equation}
  \beta_A=\frac DA\beta_D+\frac EA\beta_E=0.9091\times1.65=1.5 .
\end{equation}

\step{3}{Expected returns}
$r_A=r_U=r_{ITS}=0.05+1.5(0.10)=20\%$; $r_D=5\%$; $r_E=0.05+1.65(0.10)=21.5\%$.
MM II check: $r_E=r_U+\frac DE(r_U-r_D)=0.20+0.1(0.15)=21.5\%$. \checkmark

\answer{$\beta_A=\beta_U=\beta_{ITS}=1.5$, $\beta_D=0$. $r_A=r_U=r_{ITS}=20\%$, $r_D=5\%$,
$r_E=21.5\%$.}
\rubric{2 for stating why $\beta_{ITS}=\beta_U$ under constant $D/E$; 4 for $\beta_A=\beta_U=1.5$ from
the weights; 4 for the five returns (CAPM or MM II).}

% -----------------------------------------------------------------------------
\subsection*{(b) All-equity financing}
\statement{\QtwoB}

\step{1}{Discount rate}
The project has the risk of the comparable's unlevered assets, and with no debt $A=U=E$:
$r_A=r_E=r_U=20\%$.

\step{2}{Value and NPV}
\[
  V_U=250\times\frac{1-1.2^{-5}}{0.20}=250\times2.9906=747.65,\qquad NPV=747.65-750=-2.35 .
\]

\step{3}{Financing}
Outside investors pay at most what they receive: the whole firm is worth 747.65, so raising 750 would
require selling $750/747.65=100.31\%$ of the firm, which is impossible. No investor provides 750 for a
claim worth less than 750, and the owners would not want to undertake a negative-NPV project anyway.

\answer{$r_A=r_E=20\%$; $NPV=-2.35$; $V_U=747.65$. No: outsiders would need more than 100\%
($100.31\%$) of the firm for their 750, so the project is not financed (or undertaken) with equity
only.}
\rubric{2 for $r_A=r_E=r_U$; 4 for the annuity value and NPV; 4 for the financing argument (share
$>100\%$, so no).}

% -----------------------------------------------------------------------------
\subsection*{(c) Cost of equity at $D/E=0.5$}
\statement{\QtwoC}

With constant $D/E$, $r_{ITS}=r_U$, and MM II keeps its no-tax form:
\begin{equation}
  r_E=r_U+\frac DE(r_U-r_D)=0.20+0.5(0.20-0.05)=27.5\% .
\end{equation}
Equivalently $\beta_E=\beta_U(1+D/E)=1.5\times1.5=2.25$ and $r_E=0.05+2.25(0.10)=27.5\%$.

\answer{$r_E=27.5\%$, larger than the 20\% of (b). The same asset risk is now borne by a smaller equity
claim: debt is risk-free, so all the variability of the assets falls on equity, raising its beta from
1.5 to 2.25. The higher return is compensation for financial risk, not a higher cost of the project.}
\rubric{5 for $r_E=27.5\%$ (formula stated); 5 for ``larger'' with the financial-risk explanation.}

% -----------------------------------------------------------------------------
\subsection*{(d) NPV with debt}
\statement{\QtwoD}

\step{1}{WACC}
$D/V=0.5/1.5=1/3$ and $E/V=2/3$:
\begin{equation}
  r_{WACC}=\frac EV r_E+\frac DV r_D(1-\tau)=\tfrac23(0.275)+\tfrac13(0.05)(0.7)=0.183333+0.011667=19.5\% ,
\end{equation}
which matches the shortcut $r_{WACC}=r_U-\tau r_D\frac DV=0.20-0.30(0.05)\tfrac13=19.5\%$.

\step{2}{NPV}
\[
  V_L=250\times\frac{1-1.195^{-5}}{0.195}=250\times3.0238=755.95,\qquad NPV=755.95-750=5.95 .
\]

\answer{$NPV=+5.95$. The project's operating risk is unchanged; what changes is that interest is tax
deductible. The interest tax shields are worth more than the 2.35 shortfall of (b), and the WACC (19.5\%)
is below $r_U$ (20\%) by exactly the tax subsidy $\tau r_D D/V$. Yes, investors provide financing: the
levered firm is worth 755.95, more than the 750 they are asked to contribute.}
\rubric{4 for the WACC (either route); 3 for the NPV; 3 for the explanation (tax shield, not lower risk)
and ``yes''.}

% -----------------------------------------------------------------------------
\subsection*{(e) Market-value balance sheet at $t=0$ and the owners' stake}
\statement{\QtwoE}

\step{1}{Values}
$A_0=V_{L,0}=755.95$; $D_0=\tfrac13(755.95)=251.98$; $E_0=\tfrac23(755.95)=503.97$.

\step{2}{Who pays the 750}
Debt raises 251.98, so new shareholders must provide $750-251.98=498.02$. Paying a fair price, they
receive the fraction $498.02/503.97=98.82\%$ of the equity.

\step{3}{Owners}
The original owners keep $1-0.9882=1.18\%$ of the equity, worth $503.97-498.02=5.95$.

\answer{$D=251.98$, $E=503.97$, total assets $=755.95$. Original owners retain $1.18\%$ of the equity,
worth $5.95$: \emph{equal} to the project's NPV in (d). In competitive markets outside investors (debt
and equity) earn exactly their required return, a zero-NPV deal, so the whole NPV, tax subsidy
included, accrues to the owners who hold the project.}
\rubric{3 for $D$, $E$ and $A$; 3 for the outside share (98.82\%); 2 for the owners' share and value; 2
for ``equal'' with the competitive-markets argument.}

% -----------------------------------------------------------------------------
\subsection*{(f) Debt schedule, tax shields and $PV(ITS)$}
\statement{\QtwoF}

\step{1}{Firm value along the way}
At each date the levered firm is worth the remaining FCF discounted at the WACC,
$V_{L,t}=250\times\frac{1-1.195^{-(5-t)}}{0.195}$, and $D_t=\tfrac13V_{L,t}$.

\step{2}{Shields} $ITS_t=D_{t-1}\,r_D\,\tau=0.015\,D_{t-1}$, discounted at $r_{ITS}=r_U=20\%$.

\begin{center}
\begin{tabular}{crrrrr}
  \toprule
  $t$ & $V_{L,t}$ & $D_t$ & $ITS_t=0.015\,D_{t-1}$ & $1/1.2^t$ & $PV(ITS_t)$ \\ \midrule
  0 & 755.95 & 251.98 & -- & -- & -- \\
  1 & 653.37 & 217.79 & 3.78 & 0.8333 & 3.15 \\
  2 & 530.77 & 176.92 & 3.27 & 0.6944 & 2.27 \\
  3 & 384.27 & 128.09 & 2.65 & 0.5787 & 1.54 \\
  4 & 209.21 & 69.74 & 1.92 & 0.4823 & 0.93 \\
  5 & 0 & 0 & 1.05 & 0.4019 & 0.42 \\ \midrule
  \multicolumn{5}{r}{$PV(ITS)$ (sum of unrounded terms)} & \textbf{8.30} \\ \bottomrule
\end{tabular}
\end{center}

\step{3}{The relation}
\begin{equation}
  V_L=V_U+PV(ITS):\qquad 755.95=747.65+8.30 ,
\end{equation}
and for the project, $NPV_{(d)}=NPV_{(b)}+PV(ITS)=-2.35+8.30=5.95$ (the APV method).

\answer{$D_1=217.79$, $D_2=176.92$, $D_3=128.09$, $D_4=69.74$, $D_5=0$. Tax shields
$3.78,\ 3.27,\ 2.65,\ 1.92,\ 1.05$ at $t=1,\dots,5$, with $PV(ITS)=8.30$ at $r_U=20\%$. The levered value
equals the unlevered value of (b) plus the shields: $755.95=747.65+8.30$. WACC and APV give the same NPV
($5.95$) because both treat the shields as having the risk of $U$, as constant $D/E$ requires.}
\intuition{Debt falls as the project winds down because the firm keeps $D/E$ fixed while its value
declines; the shields shrink with it. Discounting them at $r_U$ (not $r_D$) is what makes the WACC of
(d) and the APV of (f) agree to the cent.}
\refline{Topic 2B slides (WACC vs.\ APV, $V_L=V_U+PV(ITS)$, target-leverage policies); formula sheet,
debt policies table; PS2 Exercise 2.}
\rubric{3 for the debt schedule; 3 for the shields with the lag ($D_{t-1}$); 2 for $PV(ITS)$ at $r_U$;
2 for $V_L=V_U+PV(ITS)$ linked to (b) and (d).}

\end{document}
